Importantly, practical implementations and theoretical analysis of this model rely on the metric space axioms stated in \cref{eq:metric-space-equivalence,eq:metric-space-symmetry,eq:metric-space-triangle} and not merely on Procrustes being a sensible measure of similarity as is the case of linear regression. To make this concrete, we repeated the analysis on a cohort of $K = 30$ simulated subjects whose tuning curves differed in their width (from broadly tuned to sharply tuned populations) with acuity again determined by the shape of the tuning curves using FI. Crucially, neurons across subjects additionally differed in how many other behaviorally irrelevant variables their populations encoded, as real recordings arguably do \parencite{rigotti2013,tye2024}. In other words, the overall dimensionality of the neural code differed randomly across subjects, which makes linear mapping between them one-directional: a higher-dimensional geometry can be mapped onto a lower-dimensional one, but not the reverse \parencite{elmoznino2024,muzellec2026reverse}. To quantify this, we compared Procrustes distance with a linear predictivity score ($1 - R^2$) of a cross-validated linear map from one subject's population to another's. Holding everything else fixed, kNN prediction of acuity fell from $R^2 = 0.98$ to $R^2 = -0.92$ (\textbf{Fig. } \ref{fig:synth}f, top). Symmetrizing the resulting distance matrix as $(\mbD+\mbD^\top)/2$ slightly improved the behavioral prediction, but still fell short of Procrustes ($R^2 = 0.98$ vs $R^2 = -0.13$). Using different schemes of variable encoding (linearly mixed, non-linearly mixed or segregated populations with different selectivity) slightly changed the impact on behavioral predictions (\textbf{Fig. } \ref{fig:synth}f, bottom), but predictions remained substantially worse than with a proper metric (i.e. Procrustes distance, in this case). 